\scanpage{38}
\mspage{12}

\uline{Récapitulation}

\deleted{A toute}\inserted{Soit $X$} surface conforme (sans bord), \deleted{\unsure{on peut|on peut lui},}\inserted{On peut} définir sur $X$ une métrique riemannienne canonique $Q_X$, \inserted{compatible avec str. conforme,} sauf dans les deux cas suivants :

a) $X\simeq\mathbb C$ \inserted{(i.e. $X$ droite affine sur un corps \inserted{top} $C$ isom. à $\mathbb C$)}, alors on peut définir quand même une classe de métriques riemanniennes compatibles, mod. multiplication par un scalaire. Ce sont les métriques stables par translation, qui sont aussi qu. inv. par tt autom. conforme de $X$. Elles sont \uline{plates} (donc : courbure nulle).

b) $X\simeq\mathbb P^1_{\mathbb C}$ (i.e $X$ droite projective sur un corps top $C$ \unsure{isom|$\mathbb R$-isom} à $\mathbb C$, ou encore $X$ sphère conforme). Alors il n'existe pas de métriques riemanniennes sur $X$ quasi-invariantes sous $G=\mathrm{Aut}\,X$, ni = sous $\mathrm{Aut}^+(X)$.

\deleted{Pour} Il y a pourtant sur $X$ un\deleted{e} \deleted{classe}\inserted{ens} can $R(X)$ de métriques riemanniennes (dites ``sphériques normalisées'') avec $\int_X dx=1$, qui forment une orbite sous $G$. Si $H\subset G$ est un sous-groupe, il existe un $Q\in R(X)$ stable par $H$ sss $H$ est rel cpt, et $Q$ est unique sss $X^H=\varnothing$. [Pour le cas $X^H\ne\varnothing$, cf proposition précédente qui détermine $R(x)^H$ dans tous les cas.]

% slanted
\margin{\uline{NB} $R(X)$ est un espace homogène de dim 3}
% archive sheet number: 40
